% =============================================================================
% Chapter 3 --- System Architecture
% Grounded in Context 1 §2, §7, §8 and parallel_interface_pkg.sv /
% controller_pkg.sv (for command tokens and field layout).
% =============================================================================

\chapter{System Architecture}
\label{ch:architecture}

This chapter introduces the three-module RTL stack that implements the
controller, the host-to-core packet format that every module interprets
consistently, and the repository layout in which the design and
verification code is organized. Module-internal behavior is the subject of
\Cref{ch:rtl}; this chapter deliberately stays at the system level.

\section{Top-Level Block Diagram}
\label{sec:arch-topology}

The host presents a 32-bit \emph{payload} in the same layout as the
controller's outgoing \signame{ann\_core\_word}, together with a separate
3-bit command \signame{host\_cmd}. From the host, these two signals enter
the \modname{parallel\_interface} module, which is purely combinational
and which produces a decoded bus
\{\signame{valid}, \signame{data}, \signame{address}, \signame{cmd}\}
consumed by the \modname{ann\_controller}. The controller drives the
packed output word and a 3-bit pulse-mode bus
(\signame{pulses}[2:0]) toward the ANN core, together with control lines
to the \modname{input\_buffer} that stores inference pixel data and holds
the ``expected weight'' nibble used during verify. The buffer in turn
presents one byte (\signame{buf\_data}[7:0]) for combinational reads and
eight parallel bit-serial lanes (\signame{D0}\ldots\signame{D7}) used by
the COMPUTE phase of the inference path.

\begin{figure}[htbp]
    \centering
    \includegraphics[width=0.95\linewidth]{fig_top_level_block_diagram}
    \caption[Top-level block diagram]{Top-level block diagram of the
        controller stack. The host drives \signame{host\_data} and
        \signame{host\_cmd} into the combinational
        \modname{parallel\_interface}, which produces the qualified bus
        consumed by the \modname{ann\_controller}. The controller sequences
        commands toward the ANN core over \signame{ann\_core\_word} and
        \signame{pulses}, and it also drives the \modname{input\_buffer}
        that stores expected weights and inference pixels. The
        \signame{op\_done} line is a core-side handshake sampled in the
        PROG and ERASE sub-FSMs.}
    \label{fig:top-level-block-diagram}
\end{figure}

\section{Module Roles and Responsibilities}
\label{sec:arch-modules}

\paragraph{\modname{parallel\_interface}.}
A fully combinational block. Its body contains only continuous assignments;
the clock and reset ports exist but are not referenced inside the module.
The role of the interface is two-fold. First, it slices the incoming
\signame{host\_data} into a payload byte (bits \sv{[31:24]}) and a 24-bit
one-hot \emph{tail} (bits \sv{[23:0]}) and converts the tail into a packed
16-bit \signame{address}. Second, it computes a \signame{valid} qualifier
that is asserted only when (i)~the host command is not the idle token
\sv{CMD\_HIZ} and (ii)~each of the four tail fields is strictly one-hot.
Invalid host packets therefore do not propagate past the interface, and
the downstream controller may assume that \signame{address} is legal
whenever \signame{valid} is high.

\paragraph{\modname{ann\_controller}.}
The only module in the stack that contains sequential logic. It
implements a hierarchical finite-state machine with a main FSM of nine
states (\stateval{S\_IDLE}, \stateval{S\_RESET}, \stateval{S\_PROGRAM},
\stateval{S\_VERIFY}, \stateval{S\_ERASE}, \stateval{S\_READ},
\stateval{S\_COLLECT\_DATA}, \stateval{S\_COMPUTE}, \stateval{S\_RESULT})
and three sub-FSMs for programming, verification, and erase. A shared
combinational block drives all outputs and produces next-state values;
two \sv{always\_ff} blocks register main-FSM and sub-FSM state together
with counters and captured address/data.

\paragraph{\modname{input\_buffer}.}
A synchronous 64-byte register array (\sv{BUFFER\_SIZE = 64},
\sv{BUFFER\_DATA\_WIDTH = 8}). A single mode-qualified write port loads
data from the controller during \sv{CTRL\_DATA\_LOAD} phases. Two read
paths are provided: a combinational byte-read (\signame{buf\_data}) used
by verify and read flows, and eight per-bit lanes driven from an
eight-pixel window addressed by \signame{buf\_reg\_add} and selected by
\signame{bit\_sel}[2:0], used by the INF COMPUTE phase.

\section{The 32-bit ANN Core Word}
\label{sec:arch-word}

All three modules agree on a single packet layout: the 32-bit
\signame{ann\_core\_word} carries an 8-bit payload in its high byte and a
24-bit one-hot tail in the remaining bits, where the tail encodes the
target cell as four concatenated one-hot masks
(\signame{PE}, \signame{SA}, column, row). The layout is summarized in
\Cref{tab:ann-core-word-fields} and depicted pictorially in
\Cref{fig:ann-core-word-layout}.

\begin{table}[htbp]
    \centering
    \caption{Fields of the 32-bit \signame{ann\_core\_word}.}
    \label{tab:ann-core-word-fields}
    \begin{tabularx}{\linewidth}{l l l X}
        \toprule
        \textbf{Bits} & \textbf{Width} & \textbf{Field} & \textbf{Meaning} \\
        \midrule
        \sv{[31:24]} & 8 & \signame{data\_byte} &
            Payload; for programming commands, the quantized weight nibble
            sits in bits \sv{[27:24]}. \\
        \sv{[23:20]} & 4 & \signame{PE} (one-hot) &
            Processing-element / block select. Exactly one of four bits. \\
        \sv{[19:16]} & 4 & \signame{SA} (one-hot) &
            Sub-array / sub-block select. Exactly one of four bits. \\
        \sv{[15:8]}  & 8 & Column (one-hot) &
            Column select within the sub-block. Exactly one of eight bits. \\
        \sv{[7:0]}   & 8 & Row (one-hot) &
            Row select within the sub-block. Exactly one of eight bits. \\
        \bottomrule
    \end{tabularx}
\end{table}

\begin{figure}[htbp]
    \centering
    \includegraphics[width=0.95\linewidth]{fig_ann_core_word_layout}
    \caption[ANN core word field layout]{Pictorial layout of the 32-bit
        \signame{ann\_core\_word}: an 8-bit payload byte followed by four
        one-hot fields of widths 4, 4, 8, 8. The address round-trip
        (host tail \(\rightarrow\) PI-decoded \signame{address}
        \(\rightarrow\) controller-captured \signame{address\_reg}
        \(\rightarrow\) outbound one-hot tail) is detailed in
        \Cref{ch:rtl}.}
    \label{fig:ann-core-word-layout}
\end{figure}

The 16-bit \signame{address} produced by the parallel interface is a
compact index form of the same cell identifier: it carries the packed
fields \sv{\{6'b0, blk[1:0], sb[1:0], col\_id[2:0], row\_id[2:0]\}}
(\Cref{fig:parallel-addr-layout}). The controller, when it drives the
outbound core word, reconstructs the one-hot tail by left-shifting 1 by
each index.

\begin{figure}[htbp]
    \centering
    \includegraphics[width=0.85\linewidth]{fig_parallel_addr_layout}
    \caption[Parallel-interface 16-bit address layout]{Layout of the 16-bit
        \signame{address} produced by \modname{parallel\_interface} and
        consumed by \modname{ann\_controller}. Bits \sv{[15:10]} are
        reserved and are not used by \sv{parse\_ann\_address}.}
    \label{fig:parallel-addr-layout}
\end{figure}

\section{Command Set}
\label{sec:arch-commands}

The controllable command tokens are defined in
\sv{parallel\_interface\_pkg} as an enumerated type. The five values used
in the system are listed in \Cref{tab:cmd-set}.

\begin{table}[htbp]
    \centering
    \caption{Host command set as defined in \sv{parallel\_interface\_pkg}.}
    \label{tab:cmd-set}
    \begin{tabularx}{\linewidth}{l l X}
        \toprule
        \textbf{Token} & \textbf{Encoding} & \textbf{Meaning} \\
        \midrule
        \sv{CMD\_HIZ}   & \sv{3'b000} &
            Idle; \signame{valid} is forced low at the PI regardless of
            \signame{host\_data}. \\
        \sv{CMD\_READ}  & \sv{3'b001} &
            Read the addressed cell; the controller issues a READ
            pulse train. \\
        \sv{CMD\_PROG}  & \sv{3'b010} &
            Program the addressed cell with the value in
            \signame{data\_byte}[3:0]; subject to a closed-loop
            program--verify--erase recovery policy. \\
        \sv{CMD\_ERASE} & \sv{3'b011} &
            Erase the addressed cell (host-initiated; distinct from the
            verify-initiated erase path). \\
        \sv{CMD\_INF}   & \sv{3'b100} &
            Inference; host drives nine consecutive beats of
            \signame{CMD\_INF} to collect eight pixels into the input
            buffer and trigger a COMPUTE phase. \\
        \bottomrule
    \end{tabularx}
\end{table}

The interpretation and sequencing of these tokens inside the main FSM are
described in \Cref{sec:rtl-controller-mainfsm}.

\section{Repository Organization}
\label{sec:arch-repo}

The project source, verification, and output files are organized as
follows:

\begin{itemize}
    \item \sv{source/} --- synthesizable RTL and packages, split by module:
          \sv{source/Controller/}, \sv{source/Parallel\_Interface/},
          \sv{source/Input\_Buffer/}, and a shared \sv{source/common/}
          that hosts the \sv{macros.svh} header.
    \item \sv{source/*/*\_pkg.sv} --- per-module SystemVerilog packages
          that carry functions, enumerated types, and parameter
          constants (e.g.\ \sv{controller\_pkg}).
    \item \sv{source/*/*\_rtl\_list.f} --- file lists consumed by
          \sv{vlog~-f} during RTL compilation.
    \item \sv{verif/*/tb/} and \sv{verif/*/file\_list/} --- per-module
          testbench sources and compile file lists, including both
          \emph{regular} batch benches and \emph{waves} wrappers for GUI
          debugging.
    \item \sv{verif/*/do/waves/} --- ModelSim \sv{.do} scripts that
          preload waveform groups for the wave wrappers.
    \item \sv{scripts/run\_sim.py} --- a Python driver that compiles
          modules, runs individual or bulk simulations, and performs
          housekeeping (\sv{clean}) on the \sv{work/} library.
    \item \sv{target/} --- durable simulation artifacts: pass/fail
          reports, transcripts, programming LUT images, and per-module
          log files. These are the inputs to \Cref{ch:results}.
\end{itemize}

This organization is followed strictly throughout the report: when an
RTL claim is made, a corresponding file under \sv{source/} is cited; when
a verification claim is made, a file under \sv{verif/} or \sv{target/}
is cited.
